% ==============================================================================
% Section III: Structural and Mathematical Bottlenecks
% ==============================================================================
\label{sec:bottlenecks}

Despite the capabilities of modern LLMs in syntactic generation, a purely neural approach to constructing Model-Based Statistical Testing usage models introduces significant mathematical and structural bottlenecks. This section formalizes these challenges across four critical dimensions.


% Subsection A: Stochastic Volatility and the Hallucination Dilemma

\subsection{Stochastic Volatility and the Hallucination Dilemma}
\label{subsec:hallucination}

By design, autoregressive language models predict the next token based on probabilistic vector calculations over high-dimensional latent spaces~\cite{llm_autoregressive}. This architecture is highly flexible, but it lacks the stateful execution guarantees required for formal models~\cite{llm_state_machine_frameworks}.

When translating unstructured requirements into usage models, LLMs frequently:
\begin{itemize}
    \item omit critical edge cases from the state graph,
    \item construct physically impossible state transitions, or
    \item hallucinate states that violate system invariants.
\end{itemize}

This behavior is particularly problematic for statistical usage testing~\cite{taxonomy_mbt}. Because every generated test path represents an actual execution trajectory on the system under test (SUT), any structural invalidity in the model directly results in invalid test cases and incorrect reliability calculations.

Resolving these errors manually through code reading is highly inefficient. Studies indicate that human code reviews frequently miss silent behavioral discrepancies~\cite{code_review_limits}, meaning that errors generated by the LLM often pass undetected into production test execution.


% Subsection B: The Calibration and Convex Optimization Bottleneck

\subsection{The Calibration and Convex Optimization Bottleneck}
\label{subsec:calibration}

A valid Markov chain usage model must satisfy strict mathematical constraints~\cite{jumbl_mbst_tutorial, nap_statistics}. The transition matrix $\mathbf{P}$ must be \emph{row-stochastic}, meaning that every transition probability $p_{ij}$ must satisfy:
\begin{equation}
    0 \leq p_{ij} \leq 1
    \label{eq:prob_bounds}
\end{equation}
and every row must sum to exactly one:
\begin{equation}
    \sum_{j} p_{ij} = 1, \quad \forall\, i.
    \label{eq:row_stochastic_intro}
\end{equation}

Furthermore, real-world operational profiles include complex, multi-variable constraints~\cite{boehr_constraints}. These may involve defining one transition probability as a function of another, or constraining the expected overall sequence length to a specific numerical range. Such constraints transform the calibration task into a convex optimization problem over the probability simplex.

LLMs struggle with precise numerical reasoning and constraint satisfaction~\cite{llm_numerical_reasoning}. If tasked with generating these probabilities directly, an LLM cannot guarantee that the resulting matrix satisfies the linear and convex constraints imposed by Equations~\eqref{eq:prob_bounds} and~\eqref{eq:row_stochastic_intro}. This numerical inaccuracy makes it difficult to reliably compute steady-state occupancy rates or use the model as an unbiased source of random walks for statistical testing~\cite{markov_reliability}.


% Subsection C: State-Space Explosion and Modeless Concurrency Vulnerabilities

\subsection{State-Space Explosion and Modeless Concurrency Vulnerabilities}
\label{subsec:state_explosion}

As systems scale to handle concurrent, multi-stream operations such as multi-module autonomous systems or parallel transactional checkouts where the underlying state space grows exponentially. This \emph{state-space explosion} presents a severe challenge for purely neural modeling~\cite{chatfume, protocol_gpt}.

To manage large systems, engineers have historically relied on structured notations such as the Extended Markov Modeling Language (EMML), which utilizes variables and assertions to compactly describe large usage models~\cite{emml_usage_models}. Additionally, they apply concurrency operators to combine simple, independent Markov chains into multi-stream execution patterns using specialized tools like the Test Generation Language (TGL)~\cite{jumbl_mbst_tutorial}.

Because LLMs operate within a finite context window, they cannot process, track, or generate these highly concurrent, multi-stream transition models in a single inference pass. Tasking an LLM with generating a flat, highly concurrent Markov chain usage model leads to:
\begin{enumerate}
    \item \textbf{Context fragmentation} is when the model loses track of distant state definitions as the transition matrix grows beyond the context window.
    \item \textbf{Token depletion} is a scenario where the output budget is exhausted before the full model can be articulated.
    \item \textbf{Structural coherence degradation} is when the generated graph becomes increasingly inconsistent as generation progresses, with dangling references and duplicate state identifiers.
\end{enumerate}


% Subsection D: The Semantic-Feasibility Divergence

\subsection{The Semantic-Feasibility Divergence}
\label{subsec:semantic_feasibility}

In long-horizon system validation, there is a fundamental divergence between high-level \emph{semantic progress guidance} and low-level \emph{physical feasibility verification}~\cite{neuro_symbolic_framework}. Autoregressive language models excel at semantic progress guidance: mapping out realistic pathways for users navigating an application (e.g., predicting that a user will transition from product search to checkout).

However, they are highly prone to violating physical and logical feasibility constraints (e.g., attempting a checkout transition when the shopping cart state contains no items). Without an external symbolic validation layer, a purely neural usage model will generate paths that are semantically logical but physically unexecutable on the SUT, causing high rates of false failures during automated test execution~\cite{taxonomy_mbt, markov_reliability}.

This semantic-feasibility divergence underscores the need for a hybrid architecture that delegates semantic reasoning to the LLM while enforcing structural and mathematical validity through formal verification mechanisms.
